%! The Language of Variation — Unit 1, Lesson 1.1 — Guided Notes & Practice
% THE in-class document, in the MAIN IDEAS / NOTES shape (see LESSON_SHAPE.md §1): the
% vocabulary box, then ONE two-column guidednotes table. Each row is one idea — a label on the
% left; on the right a terse definition the student completes, then a grid of short numbered
% problems with reserved work space. Problems are numbered 1..N through the whole component.
%   I DO   : the teacher fills the definition lines and works the FIRST problem of each grid.
%   WE DO  : the class works the rest of each grid, students holding the pen; the last row,
%            Guided Practice, is worked entirely together in a fresh context.
%   YOU DO : nothing on this page — ../homework/ IS the individual practice, started in the
%            last minutes of the period. No independent set, no reflection box, no objectives
%            box (the targets are on the cover).
% Vocabulary is GIVEN here, not discovered. The crux is the AMOUNTS row — a column written
% entirely in digits is not automatically quantitative. Permit numbers, parking spots, and ZIP
% codes average cleanly and mean nothing (EK VAR-1.C.1, VAR-1.C.2).
%
% THE CONTEXTS: three hiking parties (Trail / Time out), two permits with their water bottles,
% the ranger's Short / Medium / Long regrouping, the five-column digits strip, and the Season
% Pass for Guided Practice.
%
% PAGE LOCKSTEP with notes_key/: every \blank{} here becomes a short \ans{} there, every
% \termblank becomes a \termans, and every \pcell{n}{...}{H}{} gets its answer in the fourth
% argument. Heights and blank widths are sized to the answers so the two files set identically.
% The \begin{work} block is authored BYTE-IDENTICALLY in both files. A guidednotes row cannot
% break across pages, so every idea is split into two sub-rows at its prompt (a `\\` with no
% \hline, then `& ...`) so the table can break between the definition and its grid.
% NO SKETCH QUESTIONS: every display is pre-drawn and read.
\documentclass[12pt]{article}
\usepackage{apstats-article}
\usepackage{apstats-boxes}

\begin{document}

\pageheader{Unit 1, Lesson 1.1}{Guided Notes \& Practice}
% NAMESTRIP: no \namedateperiod here — the name row lives on the cover only.

\vspace{2pt}

\boxguard[16]
\begin{vocabbox}
\small
Fill in each term \textbf{as we name it} in the notes below.
\par
\termblank{Categorical variable}
\par
\termblank{Quantitative variable}
\par
\termblank{The amount test}
\par
\termblank{Grouping (binning)}
\par
\end{vocabbox}

\small
\begin{guidednotes}
% ── ROW 1: two kinds of variable ─────────────────────────────────────────────
\mainidea[Two kinds of]{Variable} &
A \textbf{variable} is any characteristic recorded for each individual; today we split them in
two. A \textbf{categorical} variable's values are category \blank{1.6cm} or group labels. A
\textbf{quantitative} variable's values are \blank{2.2cm} for a measured or counted amount.

\vspace{3pt}
Three hiking parties, two kinds of column:
\begin{center}
\renewcommand{\arraystretch}{1.1}
\begin{tabular}{@{} l l c @{}}
\rowcolor{sky}
\textbf{Party} & \textbf{Trail} & \textbf{Time out (min)} \\
\hline
Alvarez & Ridge Loop   &  95 \\
Boone   & Cedar Falls  & 140 \\
Delgado & Summit Trail & 215 \\
\hline
\end{tabular}
\end{center}
\emph{Trail} is \blank{2.6cm}: its values name which trail, and Summit Trail is not ``more''
than Ridge Loop. \emph{Time out} is \blank{2.6cm}: 215 minutes really is more than 95, and the
difference, 120 minutes, means something.
\\
& \notesprompt{The ranger records more about each party. Classify each: C or Q.}
\begin{probgrid}{|Y|Y|Y|}
\pcell{1}{Group size (number of hikers)}{0.9cm}{} &
\pcell{2}{Dog along (Yes / No)}{0.9cm}{} &
\pcell{3}{Weather at the trailhead}{0.9cm}{} \\ \hline
\pcell{4}{Pack weight, in pounds}{1.1cm}{} &
\pcell{5}{Trail difficulty (Easy / Moderate / Hard)}{1.1cm}{} &
\pcell{6}{Miles hiked}{1.1cm}{} \\ \hline
\end{probgrid}
\\ \hline
% ── ROW 2: the amount test ───────────────────────────────────────────────────
\mainidea[The]{Amount test} &
Do not ask ``does it have digits.'' Ask: \textbf{is the value an \emph{amount} of something ---
how much, or how many?} If yes, the variable is \blank{2.8cm}. If the value only
\emph{identifies} or \emph{labels}, it is \blank{2.6cm}.

\vspace{3pt}
A second way to run the same test: \textbf{does averaging the values produce something real?}
If the average is \blank{2.2cm}, the variable is categorical no matter how it is written.

\\
& Two hiking parties hold permits 2041 and 6318 and carry 6 and 2 water bottles. Both columns are
written entirely in digits. Both averages compute without complaint:

\begin{work}
\text{average permit \#} &= \frac{2041 + 6318}{2} = 4179.5 \\[4pt]
\text{average water bottles} &= \frac{6 + 2}{2} = 4
\end{work}
\notesprompt{Run the test on both averages.}
\begin{probgrid}{|Y|Y|Y|}
\pcell{7}{Four water bottles --- what real thing does that number describe?}{1.5cm}{} &
\pcell{8}{Hand me the party whose permit is 4179.5.}{1.5cm}{} &
\pcell{9}{Which column passes the amount test, and what is it an amount of?}{1.5cm}{} \\ \hline
\end{probgrid}
\\ \hline
% ── ROW 3: grouping runs one way ─────────────────────────────────────────────
\mainidea[Grouping is]{One way only} &
A ranger stops recording \emph{Time out} in minutes and records \emph{Short} / \emph{Medium} /
\emph{Long} instead. Same hikes, same afternoon:

\vspace{2pt}
\centerline{\emph{Time out} \ 95, 140, 215 \ minutes \quad$\longrightarrow$\quad
\emph{Length} \ Short, Medium, Long}

\vspace{2pt}
The new column is \blank{2.6cm}. Replacing measured amounts with group labels is called
\textbf{grouping} or \textbf{binning}. Once 95, 140, and 215 become labels, the mean time on the
trail can no longer be \blank{2.1cm} from the column. Grouping runs \blank{1.2cm} way only:
you can always bin amounts into labels, but you can never \blank{1.8cm} the amounts from the
labels.
\\
& \notesprompt{Use the two columns.}
\begin{probgrid}{|Y|Y|}
\pcell{10}{Which column can be averaged? Compute that average.}{1.3cm}{} &
\pcell{11}{Boone's row now reads \emph{Medium}. Can the ranger get 140 back from it? Why?}{1.3cm}{} \\ \hline
\pcell{12}{Name one thing the ranger can no longer compute from the new column.}{1.2cm}{} &
\pcell{13}{A second ranger records \emph{Under 2 hours} / \emph{2 hours or more}. Which kind of
variable is that column?}{1.2cm}{} \\ \hline
\end{probgrid}
\\ \hline
% ── ROW 4: digits are not amounts — THE CRUX ─────────────────────────────────
\mainidea[Digits are not]{Amounts} &
A column written entirely in digits is \blank{1.2cm} automatically quantitative. Permit numbers
and ZIP codes average without complaint and mean \blank{1.6cm}: their digits \blank{1.7cm} an
individual, they do not measure anything.

\vspace{3pt}
Counting a categorical variable also produces numbers. ``62\% of parties hiked with a dog'' is a
numerical \blank{1.9cm} of the categorical variable \emph{Dog along}. The \emph{summary} is
numeric; the \emph{variable} is still \blank{2.4cm}.
\\
& \notesprompt{Classify each. Write C or Q.}
\begin{probgrid}{|Y|Y|Y|Y|Y|}
\pcell{14}{Permit \#}{0.9cm}{} &
\pcell{15}{Parking spot}{0.9cm}{} &
\pcell{16}{Home ZIP}{0.9cm}{} &
\pcell{17}{Water bottles}{0.9cm}{} &
\pcell{18}{Group size}{0.9cm}{} \\ \hline
\end{probgrid}
\notesprompt{Then say why.}
\begin{probgrid}{|Y|Y|}
\pcell{19}{Every entry in the \emph{Dog along} column is Yes or No. Where did the 62 come from
--- and is it a value of the variable?}{1.4cm}{} &
\pcell{20}{State the amount test in one sentence --- without using the word
\emph{digits}.}{1.4cm}{} \\ \hline
\end{probgrid}
\\ \hline
% ── ROW 5: guided practice — WE DO, together, fresh context ──────────────────
\mainidea[Guided practice]{The Season Pass} &
A ski area records five things about each of its 3{,}400 season-pass holders: \textbf{pass
number}, \textbf{age} (years), \textbf{pass tier} (Gold / Silver / Bronze), \textbf{miles
traveled} to the mountain, and \textbf{locker number}.

\vspace{4pt}
\textbf{21.} The individuals are \blank{3.6cm}.

\vspace{4pt}
\textbf{22.} Label each variable \textbf{C} or \textbf{Q}:

\vspace{2pt}
\renewcommand{\arraystretch}{1.3}
\begin{tabularx}{\linewidth}{@{} Y Y Y Y Y @{}}
\textbf{Pass \#} & \textbf{Age} & \textbf{Pass tier} & \textbf{Miles traveled} &
\textbf{Locker \#} \\
\blank{1.4cm} & \blank{1.4cm} & \blank{1.4cm} & \blank{1.4cm} & \blank{1.4cm} \\
\end{tabularx}
\\
& \notesprompt{We work these together.}
\begin{probgrid}{|Y|Y|Y|}
\pcell{23}{\emph{Miles traveled} and \emph{Locker \#} are both digits. Say the difference in
one sentence.}{1.6cm}{} &
\pcell{24}{The resort stops recording age in years and records \textbf{age group} instead
(Under 18 / 18--34 / 35--54 / 55+). Which kind is the variable now, and what can the resort no
longer compute?}{1.6cm}{} &
\pcell{25}{A report states: \textbf{``The average pass number is 5{,}120.''} Why does that
number tell you nothing about the pass holders?}{1.6cm}{} \\ \hline
\end{probgrid}
\\ \hline
\end{guidednotes}

% ── THE NOTES END AT GUIDED PRACTICE ─────────────────────────────────────────
% There is deliberately no "you do" here: ../homework/ is the individual practice,
% started in class in the last 8 minutes while the teacher circulates.

\end{document}
